\documentclass{amsart}
% For dealing with references we use the comment environment
\usepackage{verbatim}
\newenvironment{reference}{\comment}{\endcomment}
%\newenvironment{reference}{}{}
\newenvironment{slogan}{\comment}{\endcomment}
\newenvironment{history}{\comment}{\endcomment}

% For commutative diagrams we use Xy-pic
\usepackage[all]{xy}

% We use 2cell for 2-commutative diagrams.
\xyoption{2cell}
\UseAllTwocells

% We use multicol for the list of chapters between chapters
\usepackage{multicol}

% This is generally recommended for better output
\usepackage{lmodern}
\usepackage[T1]{fontenc}

% For cross-file-references

% Package for hypertext links:
\usepackage{hyperref}

% For any local file, say "hello.tex" you want to link to please

% Theorem environments.
%
\theoremstyle{plain}
\newtheorem{theorem}[subsection]{Theorem}
\newtheorem{proposition}[subsection]{Proposition}
\newtheorem{lemma}[subsection]{Lemma}

\theoremstyle{definition}
\newtheorem{definition}[subsection]{Definition}
\newtheorem{example}[subsection]{Example}
\newtheorem{exercise}[subsection]{Exercise}
\newtheorem{situation}[subsection]{Situation}

\theoremstyle{remark}
\newtheorem{remark}[subsection]{Remark}
\newtheorem{remarks}[subsection]{Remarks}

\numberwithin{equation}{subsection}

% Macros
%
\def\lim{\mathop{\mathrm{lim}}\nolimits}
\def\colim{\mathop{\mathrm{colim}}\nolimits}
\def\Spec{\mathop{\mathrm{Spec}}}
\def\Hom{\mathop{\mathrm{Hom}}\nolimits}
\def\Ext{\mathop{\mathrm{Ext}}\nolimits}
\def\SheafHom{\mathop{\mathcal{H}\!\mathit{om}}\nolimits}
\def\SheafExt{\mathop{\mathcal{E}\!\mathit{xt}}\nolimits}
\def\Sch{\mathit{Sch}}
\def\Mor{\mathop{\mathrm{Mor}}\nolimits}
\def\Ob{\mathop{\mathrm{Ob}}\nolimits}
\def\Sh{\mathop{\mathit{Sh}}\nolimits}
\def\NL{\mathop{N\!L}\nolimits}
\def\CH{\mathop{\mathrm{CH}}\nolimits}
\def\proetale{{pro\text{-}\acute{e}tale}}
\def\etale{{\acute{e}tale}}
\def\QCoh{\mathit{QCoh}}
\def\Ker{\mathop{\mathrm{Ker}}}
\def\Im{\mathop{\mathrm{Im}}}
\def\Coker{\mathop{\mathrm{Coker}}}
\def\Coim{\mathop{\mathrm{Coim}}}

% Boxtimes
%
\DeclareMathSymbol{\boxtimes}{\mathbin}{AMSa}{"02}

%
% Macros for moduli stacks/spaces
%
\def\QCohstack{\mathcal{QC}\!\mathit{oh}}
\def\Cohstack{\mathcal{C}\!\mathit{oh}}
\def\Spacesstack{\mathcal{S}\!\mathit{paces}}
\def\Quotfunctor{\mathrm{Quot}}
\def\Hilbfunctor{\mathrm{Hilb}}
\def\Curvesstack{\mathcal{C}\!\mathit{urves}}
\def\Polarizedstack{\mathcal{P}\!\mathit{olarized}}
\def\Complexesstack{\mathcal{C}\!\mathit{omplexes}}
% \Pic is the operator that assigns to X its picard group, usage \Pic(X)
% \Picardstack_{X/B} denotes the Picard stack of X over B
% \Picardfunctor_{X/B} denotes the Picard functor of X over B
\def\Pic{\mathop{\mathrm{Pic}}\nolimits}
\def\Picardstack{\mathcal{P}\!\mathit{ic}}
\def\Picardfunctor{\mathrm{Pic}}
\def\Deformationcategory{\mathcal{D}\!\mathit{ef}}

\setlength{\emergencystretch}{3em}
\begin{document}
\title[Stacks Project, Tag 0200]{016 --- Chow lemma for separated finite-type morphisms over a Noetherian base}
\maketitle
\noindent Copyright (C) 2005--2025 Johan de Jong. Stacks Project authors. Source: \href{https://stacks.math.columbia.edu/tag/0200}{Tag 0200}.

\noindent Editorial difficulty note (not author text). Difficulty H2: The proof constructs a quasi-projective proper modification of a separated finite-type scheme over a Noetherian base using closures of diagonal graph maps. Scheme-theoretic density permits gluing the local projections and proves exact inverse-image identities; these verify properness, surjectivity and a dense-open isomorphism. This is a substantive geometric construction beyond routine qualifying algebra..

\noindent Editorial provenance note (not author text). History: excerpt from revision \texttt{a0444\allowbreak 6e57e\allowbreak c1fbc\allowbreak 252a8\allowbreak 71afc\allowbreak ec775\allowbreak 2fb28\allowbreak 07b14}, coherent.tex, lines 4922--5072. Reference presentation was adapted; original-author context and core support blocks were retained or added. The mathematical wording is unchanged. Licensed under GNU FDL 1.2 or later; no invariant sections or cover texts. See ../licenses/COPYING. Context blocks reproduce author definitions and supporting results; they are not additional counted problems.

\begin{definition}
\label{morphisms-definition-scheme-theoretic-image}
Let $f : X \to Y$ be a morphism of schemes. The {\it scheme theoretic image}
of $f$ is the smallest closed subscheme $Z \subset Y$ through which $f$
factors, see Lemma \href{https://stacks.math.columbia.edu/tag/01R6}{lemma scheme theoretic image (Tag 01R6)} above.
\end{definition}

\begin{definition}
\label{morphisms-definition-scheme-theoretically-dense}
Let $X$ be a scheme. Let $U \subset X$ be an open subscheme.
\begin{enumerate}
\item The scheme theoretic image of the morphism $U \to X$
is called the {\it scheme theoretic closure of $U$ in $X$}.
\item We say $U$ is {\it scheme theoretically dense in $X$}
if for every open $V \subset X$ the scheme theoretic closure
of $U \cap V$ in $V$ is equal to $V$.
\end{enumerate}
\end{definition}

\begin{lemma}
\label{lemma-chow-Noetherian}
\begin{reference}
\href{https://stacks.math.columbia.edu/bibliography}{Bibliography: EGA (II Theorem 5.6.1(a))}
\end{reference}
Let $S$ be a Noetherian scheme.
Let $f : X \to S$ be a separated morphism of finite type.
Then there exist an $n \geq 0$ and a diagram
$$
\xymatrix{
X \ar[rd] & X' \ar[d] \ar[l]^\pi \ar[r] & \mathbf{P}^n_S \ar[dl] \\
& S &
}
$$
where $X' \to \mathbf{P}^n_S$ is an immersion, and
$\pi : X' \to X$ is proper and surjective. Moreover, we may
arrange it such that there exists a dense open subscheme
$U \subset X$ such that $\pi^{-1}(U) \to U$ is an isomorphism.
\end{lemma}

\begin{proof}
All of the schemes we will encounter during the rest of the proof
are going to be of finite type over the Noetherian scheme $S$ and
hence Noetherian
(see Morphisms, Lemma \href{https://stacks.math.columbia.edu/tag/01T6}{lemma finite type noetherian (Tag 01T6)}).
All morphisms between them will automatically be quasi-compact, locally of
finite type and quasi-separated, see
Morphisms, Lemma \href{https://stacks.math.columbia.edu/tag/01T8}{lemma permanence finite type (Tag 01T8)} and
Properties,
Lemmas \href{https://stacks.math.columbia.edu/tag/01OY}{lemma locally Noetherian quasi separated (Tag 01OY)} and
\href{https://stacks.math.columbia.edu/tag/01P0}{lemma morphism Noetherian schemes quasi compact (Tag 01P0)}.

\medskip\noindent
The scheme $X$ has only finitely many irreducible components
(Properties, Lemma \href{https://stacks.math.columbia.edu/tag/0BA8}{lemma Noetherian irreducible components (Tag 0BA8)}).
Say $X = X_1 \cup \ldots \cup X_r$ is the decomposition
of $X$ into irreducible components.
Let $\eta_i \in X_i$ be the generic point.
For every point $x \in X$ there exists an affine open
$U_x \subset X$ which contains $x$ and each of the generic
points $\eta_i$. See
Properties, Lemma \href{https://stacks.math.columbia.edu/tag/01ZX}{lemma point and maximal points affine (Tag 01ZX)}.
Since $X$ is quasi-compact, we can find a finite affine open
covering $X = U_1 \cup \ldots \cup U_m$ such that
each $U_i$ contains $\eta_1, \ldots, \eta_r$.
In particular we conclude that the open
$U = U_1 \cap \ldots \cap U_m \subset X$ is
a dense open. This and the fact that the $U_i$ are affine opens
covering $X$ are all that we will use below.

\medskip\noindent
Let $X^* \subset X$ be the scheme theoretic closure of $U \to X$, see
Morphisms, Definition \ref{morphisms-definition-scheme-theoretic-image}.
Let $U_i^* = X^* \cap U_i$. Note that $U_i^*$ is a closed subscheme
of $U_i$. Hence $U_i^*$ is affine. Since $U$ is dense in $X$ the
morphism $X^* \to X$ is a surjective closed immersion. It is an
isomorphism over $U$. Hence we may replace $X$ by $X^*$ and
$U_i$ by $U_i^*$ and assume that $U$ is scheme theoretically dense
in $X$, see
Morphisms, Definition \ref{morphisms-definition-scheme-theoretically-dense}.

\medskip\noindent
By Morphisms, Lemma \href{https://stacks.math.columbia.edu/tag/01VS}{lemma quasi projective finite type over S (Tag 01VS)}
we can find an immersion $j_i : U_i \to \mathbf{P}_S^{n_i}$
for each $i$. By
Morphisms, Lemma \href{https://stacks.math.columbia.edu/tag/01RG}{lemma quasi compact immersion (Tag 01RG)} we can find
closed subschemes $Z_i \subset \mathbf{P}_S^{n_i}$
such that $j_i : U_i \to Z_i$ is a scheme theoretically
dense open immersion. Note that $Z_i \to S$ is proper, see
Morphisms, Lemma \href{https://stacks.math.columbia.edu/tag/01WC}{lemma locally projective proper (Tag 01WC)}.
Consider the morphism
$$
j = (j_1|_U, \ldots, j_m|_U) : U \longrightarrow
\mathbf{P}_S^{n_1} \times_S \ldots \times_S \mathbf{P}_S^{n_m}.
$$
By the lemma cited above we can find a closed subscheme
$Z$ of $\mathbf{P}_S^{n_1} \times_S \ldots \times_S \mathbf{P}_S^{n_m}$
such that $j : U \to Z$ is an open immersion and such that $U$
is scheme theoretically dense in $Z$. The morphism $Z \to S$
is proper. Consider the $i$th projection
$$
\text{pr}_i|_Z : Z \longrightarrow \mathbf{P}^{n_i}_S.
$$
This morphism factors through $Z_i$ (see Morphisms,
Lemma \href{https://stacks.math.columbia.edu/tag/01R9}{lemma factor factor (Tag 01R9)}). Denote $p_i : Z \to Z_i$
the induced morphism. This is a proper morphism, see
Morphisms, Lemma \href{https://stacks.math.columbia.edu/tag/01W6}{lemma image proper scheme closed (Tag 01W6)}
for example. At this point we have that
$U \subset U_i \subset Z_i$ are scheme theoretically
dense open immersions. Moreover, we can think of $Z$ as the
scheme theoretic image of the ``diagonal'' morphism
$U \to Z_1 \times_S \ldots \times_S Z_m$.

\medskip\noindent
Set $V_i = p_i^{-1}(U_i)$. Note that $p_i|_{V_i} : V_i \to U_i$ is proper.
Set $X' = V_1 \cup \ldots \cup V_m$. By construction $X'$ has an immersion
into the scheme
$\mathbf{P}^{n_1}_S \times_S \ldots \times_S \mathbf{P}^{n_m}_S$.
Thus by the Segre embedding (see
Constructions, Lemma \href{https://stacks.math.columbia.edu/tag/01WD}{lemma segre embedding (Tag 01WD)})
we see that $X'$ has
an immersion into a projective space over $S$.

\medskip\noindent
We claim that the morphisms $p_i|_{V_i}: V_i \to U_i$ glue to a morphism
$X' \to X$. Namely, it is clear that $p_i|_U$ is the identity map
from $U$ to $U$. Since $U \subset X'$ is scheme theoretically
dense by construction, it is also scheme theoretically dense
in the open subscheme $V_i \cap V_j$. Thus we see that
$p_i|_{V_i \cap V_j} = p_j|_{V_i \cap V_j}$ as morphisms into the
separated $S$-scheme $X$, see
Morphisms, Lemma \href{https://stacks.math.columbia.edu/tag/01RH}{lemma equality of morphisms (Tag 01RH)}.
We denote the resulting morphism $\pi : X' \to X$.

\medskip\noindent
We claim that $\pi^{-1}(U_i) = V_i$.
Since $\pi|_{V_i} = p_i|_{V_i}$ it follows that
$V_i \subset \pi^{-1}(U_i)$. Consider the diagram
$$
\xymatrix{
V_i \ar[r] \ar[rd]_{p_i|_{V_i}} & \pi^{-1}(U_i) \ar[d] \\
& U_i
}
$$
Since $V_i \to U_i$ is proper we see that the image of
the horizontal arrow is closed, see
Morphisms, Lemma \href{https://stacks.math.columbia.edu/tag/01W6}{lemma image proper scheme closed (Tag 01W6)}.
Since $V_i \subset \pi^{-1}(U_i)$ is scheme
theoretically dense (as it contains $U$)
we conclude that $V_i = \pi^{-1}(U_i)$ as claimed.

\medskip\noindent
This shows that $\pi^{-1}(U_i) \to U_i$ is identified with the proper
morphism $p_i|_{V_i} : V_i \to U_i$. Hence we see that $X$ has a
finite affine covering $X = \bigcup U_i$ such that the restriction
of $\pi$ is proper on each member of the covering.
Thus by Morphisms, Lemma \href{https://stacks.math.columbia.edu/tag/01W2}{lemma proper local on the base (Tag 01W2)}
we see that $\pi$ is proper.

\medskip\noindent
Finally we have to show that $\pi^{-1}(U) = U$. To see this we argue in the
same way as above using the diagram
$$
\xymatrix{
U \ar[r] \ar[rd] & \pi^{-1}(U) \ar[d] \\
& U
}
$$
and using that $\text{id}_U : U \to U$ is proper and that
$U$ is scheme theoretically dense in $\pi^{-1}(U)$.
\end{proof}
\end{document}
